
\section{On Withdrawal from an “Association"}

Editorial article in “Folkebladet” for February 23, 1898.

We printed Pastor A. Wright’s piece on “going out” in the last number of “Folkebladet.” And we would gladly ask many to read it again. It contains much that is good, if it is applied to withdrawal from a congregation; at the same time that it is exceedingly weak when it is applied to withdrawal from an “association.”

We will also make a few remarks upon the piece, which we hope will bring Christian people onto the right track in this matter.

Concerning the reasons why one ought not to go out of “our association,” Pastor Wright says among other things:

“1. God wills that His Church should be one, not only in spirit but also in body. John 17. 1 Cor. Frailties in persons within the association do not set aside this will of God.”

It appears that Pastor Wright means that Jesus in John 17, the so-called High-Priestly Prayer, prayed for and concerning an outward organization of the Church, so that this prayer is fulfilled by all congregations joining themselves to one and the same outward organization, and then preferably the United Church.

Here, then, the question arises:
Does the formation of the so-called “associations” promote or hinder the Church’s inward and outward unity?

We know from the New Testament that the Church’s unity is not hindered by the many congregations. Jesus and His Apostles have not objected to the number of congregations increasing in a lawful way.

But if association is formed after association, and association against association, and each of them says of itself that it is the one true Church, does not this hinder the Church’s unity?

Can it be regarded as absolutely certain that the formation of these ruling church bodies really belonged to the plan of Jesus and the Apostles?

But if it should be so that the freedom and independence of congregations, without any ruling “association” over them, is the original order, the one most serviceable to the Church’s unity, then it surely must be asked whether those who tighten the bond of the association most and make it such a matter of conscience are also serving the Church’s unity best.

For we all know that congregations are split and torn asunder for the sake of the so-called associations; is this then also a promotion of the Church’s unity?

And further Pastor Wright says:
“4. To go out of the congregation or association I have the right only” etc.

To this we will only remark that this view, according to which congregation and association are placed absolutely side by side, seems to us either to set the congregation too low or the association too high.

“Folkebladet” does not wish to pronounce upon the many reasons that speak for withdrawal from the United Church. We are sure that the congregations will find out where the shoe pinches. If there are any among them who feel that they cannot with a good conscience stand with the “association,” then they know well enough what they are to do. “Folkebladet” has no concern for the purely outward side of the matter.

Georg Sverdrup
